\documentclass[10pt]{beamer}
\usetheme{metropolis} 
\usecolortheme{default}

\usepackage{booktabs}
\usepackage{amsmath}
\usepackage{graphicx}
\usepackage{multirow}

\title{Welcome to Public Choice}
\subtitle{ECO1028 - Politics Without Romance}
\author{Beatriz Gietner\\
\texttt{beatriz.gietner@dcu.ie}}
\date{\today}

\begin{document}

\maketitle

\begin{frame}{Before We Begin...}
\textbf{Let's get to know each other!}

\vspace{0.5cm}
Go around the room:
\begin{itemize}
\item Your name
\item Your programme (Politics/Law/Economics/Other)
\item One political puzzle that interests you
\item What you hope to learn from this module
\end{itemize}

\vspace{0.5cm}
\small{(We'll take about 15--20 minutes for this)}
\end{frame}

\begin{frame}{About Me}
\textbf{Background:}
\begin{itemize}
\item Applied microeconomist working on education, labour and econometrics
\item \textbf{Not Irish!} (But living here now)
\item This means: \textit{You} know more about Irish politics than I do
\item Therefore... you'll teach me about Ireland, I'll teach you the economic theory
\end{itemize}

\vspace{0.3cm}

\textbf{Office Hours:}
\begin{itemize}
\item \textbf{Fridays:} 11:30--12:30 via Zoom
\begin{itemize}
\item[\textbullet] \small{\url{https://dcu-ie.zoom.us/j/86881081383?pwd=9Lz5U5qQL9lCbS3nTCxgp15lewudVs.1}}
\item[\textbullet] \small{Meeting ID: 868 8108 1383 | Passcode: 676403}
\end{itemize}
\item \textbf{Wednesdays:} 11:00--12:00 in person at Prof Robert's office (Q251)
\end{itemize}

\vspace{0.2cm}

\textbf{Email:} \texttt{beatriz.gietner@dcu.ie} (24--48 hour response time)
\end{frame}

\begin{frame}{Lecture Outline}
\tableofcontents
\end{frame}

\section{What is Public Choice?}

\begin{frame}{A Provocative Question}
\begin{center}
\Large Why do governments fail?
\end{center}

\pause
\vspace{1cm}

Economics 101 tells us markets can fail:
\begin{itemize}
\item Monopoly
\item Externalities
\item Public goods
\item Information asymmetries
\end{itemize}

\pause
\textbf{But what about government failure?}
\end{frame}

\begin{frame}{Defining the Field}
\begin{quote}
"Public choice is politics without romance... The wishful thinking it displaced presumes that participants in the political sphere aspire to promote the common good."
\end{quote}
\begin{flushright}
- James Buchanan (Nobel Laureate, 1986)
\end{flushright}

\pause
\vspace{0.5cm}
\begin{itemize}
\item \textbf{The Core Insight:} we apply economic tools (rationality, incentives, constraints) to political behaviour
\item Treat politicians, voters, and bureaucrats as \textit{humans} responding to incentives, not saints
\item If self-interest explains market behaviour, why not political behaviour?
\end{itemize}
\end{frame}

\begin{frame}{The Paradigm Shift}
\begin{columns}
\begin{column}{0.48\textwidth}
\textbf{Traditional View}
\begin{itemize}
\item Government pursues "public interest"
\item Politicians are benevolent social planners
\item Government fixes market failures
\item Focus on policy outcomes
\end{itemize}
\end{column}
\pause
\begin{column}{0.48\textwidth}
\textbf{Public Choice View}
\begin{itemize}
\item Actors pursue private interests
\item Politicians maximise votes/power
\item Government has its own failures
\item Focus on political \textit{processes}
\end{itemize}
\end{column}
\end{columns}

\vspace{0.5cm}
\pause
\textbf{Key insight:} The same person who behaves selfishly in the market doesn't suddenly become altruistic in the voting booth.
\end{frame}

\begin{frame}{An Example: Why Do We Have Tariffs?}
\textbf{Traditional explanation:}
\begin{itemize}
\item Protect infant industries
\item National security
\item Economic development
\end{itemize}

\pause

\textbf{Public Choice explanation:}
\begin{itemize}
\item Concentrated benefits (protected industry) vs.\ dispersed costs (consumers)
\item Protected industry lobbies intensely; consumers remain rationally ignorant
\item Politicians respond to organised interests with votes and campaign contributions
\item Result: inefficient policy persists
\end{itemize}

\pause
\vspace{0.4cm}
\textbf{Takeaway:} Public Choice shifts the question from \textit{what should be done?} to \textit{what incentives drive the actors?}
\end{frame}

\section{A Brief History of Public Choice}

\begin{frame}{The Origins: 1940s--1960s}
\textbf{The Founders:}
\begin{itemize}
\item \textbf{(Philosophical precursor) Knut Wicksell (1896):} unanimity and "just taxation" as inspiration for constitutional political economy
\item \textbf{(Precursor) Duncan Black (1948):} "On the Rationale of Group Decision-making" -- median voter theorem
\item \textbf{Kenneth Arrow (1951):} Impossibility theorem -- no perfect voting system (Nobel 1972)
\item \textbf{Anthony Downs (1957):} \textit{An Economic Theory of Democracy} -- rational voter model
\item \textbf{Mancur Olson (1965):} \textit{The Logic of Collective Action} -- free-rider problem
\end{itemize}


\textbf{The breakthrough:} Applying economic reasoning to political behaviour


\small{Before this, political science and economics were separate worlds.}
\end{frame}

\begin{frame}{The Golden Age: 1960s--1980s}
\textbf{Virginia School (Buchanan \& Tullock):}
\begin{itemize}
\item \textit{The Calculus of Consent} (1962) -- constitutional political economy
\item Buchanan wins Nobel Prize (1986)
\item Focus: How should political institutions be designed?
\item More normative, constitutional emphasis
\end{itemize}


\textbf{Other key contributions:}
\begin{itemize}
\item \textbf{Gordon Tullock (1967):} Rent-seeking theory
\item \textbf{William Niskanen (1971):} Budget-maximising bureaucrat
\item \textbf{George Stigler (1971):} Theory of regulation and capture
\end{itemize}


\textbf{The main insight:} Government agents respond to incentives just like market actors.
\end{frame}

\begin{frame}{Modern Developments: 1990s--Present}
\textbf{Three major trends:}


\textbf{1. Empirical Revolution}
\begin{itemize}
\item Testing theories with real data (elections, policy outcomes)
\item Natural experiments and quasi-experimental methods
\item Example: Do political business cycles actually exist?
\end{itemize}


\textbf{2. Behavioural Public Choice}
\begin{itemize}
\item Incorporating insights from psychology
\item Bounded rationality, cognitive biases
\item Example: Nudges and choice architecture (Thaler \& Sunstein)
\end{itemize}


\textbf{3. Institutions \& Development (Modern Political Economy)}
\begin{itemize}
\item Why are some countries rich and others poor?
\item Role of political institutions in economic development
\item Example: Acemoglu \& Robinson's work on extractive vs.\ inclusive institutions
\end{itemize}
\end{frame}

\begin{frame}{Key Debates in Public Choice}
\textbf{1. Market failure vs.\ government failure}
\begin{itemize}
\item Traditional economics: Government fixes market failures
\item Public Choice: But government has its own failures
\item Which is worse? Context-dependent
\end{itemize}


\textbf{2. Positive vs.\ normative analysis}
\begin{itemize}
\item Positive: How do political systems actually work?
\item Normative: How should they work? Constitutional design
\item Both matter
\end{itemize}


\textbf{3. Rational choice vs.\ behavioural approaches}
\begin{itemize}
\item Are people rational optimisers or boundedly rational satisficers?
\item When do biases matter for political outcomes?
\item Ongoing debate -- we'll explore both sides
\end{itemize}
\end{frame}

\begin{frame}{Public Choice Today}
\textbf{Where is the field now?}
\begin{itemize}
\item Major journals: \textit{Public Choice}, \textit{Journal of Public Economics}, political science journals
\item Interdisciplinary: Economics, political science, law, sociology
\item Policy relevant: Institutional design, constitutional reform, development policy
\end{itemize}


\textbf{Current hot topics:}
\begin{itemize}
\item Populism and polarisation
\item Social media and political behaviour
\item Climate change and collective action
\item Democratic backsliding
\item Corruption and governance in developing countries
\end{itemize}


\textbf{The toolkit is more relevant than ever.}
\end{frame}

\section{Foundational Principles}

\begin{frame}{The Three Pillars of Public Choice}
\begin{enumerate}
\item \textbf{Methodological Individualism}
\begin{itemize}
\item Only individuals choose and act
\item "The State," "The Party," "The People" are collections of individuals
\item We analyse individual decision-makers rather than abstract entities
\end{itemize}

\pause
\item \textbf{Rational Choice}
\begin{itemize}
\item People maximise their own utility subject to constraints
\item True in markets \textit{and} in politics
\item Doesn't mean people are selfish... they are just consistent
\end{itemize}

\pause
\item \textbf{Politics as Exchange}
\begin{itemize}
\item Markets exchange goods and services
\item Politics exchanges votes, taxes, and policies
\item Both are mechanisms for coordination under scarcity
\end{itemize}
\end{enumerate}
\end{frame}

\begin{frame}{But Wait... Are People Really Rational?}
\textbf{Common objections:}
\begin{itemize}
\item "People vote against their own interests!"
\item "People are emotional and biased!"
\item "Politicians care about legacy, not just winning!"
\end{itemize}

\pause
\vspace{0.5cm}
\textbf{Our approach:}
\begin{itemize}
\item We'll start with the rational choice model (Weeks 2--6, 8)
\item Then we'll challenge it with behavioural economics (Week 9)
\item We start with rational choice \textbf{because} it is a clear benchmark
\item Once we understand its predictions, we can see where and why reality diverges
\end{itemize}

\pause
\small{As Milton Friedman argued: "the only relevant test of the validity of a hypothesis is comparison of its predictions with experience" (1953). Do rational choice models predict behaviour? We'll see.}
\end{frame}
\section{Course Structure}

\begin{frame}{Our Journey: 12 Weeks}
\textbf{Weeks 1--10:} Theory + Irish Applications
\begin{itemize}
\item Each week: New public choice topic
\item Student groups present Irish context (Weeks 3--6, 8--10)
\item Week 7 = Reading Week (no class)
\end{itemize}

\textbf{Weeks 11--12:} Your presentations (Project 1)

\vspace{0.5cm}
\pause
\textbf{The Flow:}
\begin{enumerate}
\item Demand side: How do voters behave?
\item Supply side: How do politicians and bureaucrats respond?
\item Interest groups: Who gets policy and why?
\item Structures: How do institutions shape outcomes?
\item Challenges: Is our theory right?
\item Synthesis: What determines good governance across different regimes?
\end{enumerate}
\end{frame}


\begin{frame}{Week-by-Week Schedule (1/2)}
\small
\begin{tabular}{p{0.9cm} p{3.1cm} p{5.8cm}}
\toprule
\textbf{Week} & \textbf{Topic (t)} & \textbf{Irish Context Presentation (applies Week $t-1$)} \\
\midrule
1 & Introduction & - \\
2 & Paradox of Voting & - \\
3 & Median Voter Theorem & Voting behaviour (turnout, participation, abstention) \\
4 & Political Competition and Macroeconomic Performance & Party positioning; PR-STV; policy convergence \\
5 & Bureaucracy & Political budget cycles; opportunistic spending; pre-election budgets \\
6 & Federalism \& Decentralisation & HSE; civil service; agency problems; centralisation \\
\bottomrule
\end{tabular}
\end{frame}

\begin{frame}{Week-by-Week Schedule (2/2)}
\small
\begin{tabular}{p{0.9cm} p{3.1cm} p{5.8cm}}
\toprule
\textbf{Week} & \textbf{Topic (t)} & \textbf{Irish Context Presentation (applies Week $t-1$)} \\
\midrule
\textbf{7} & \textbf{READING WEEK} & \textbf{No class - prepare Project 1 presentations} \\
\addlinespace
8 & Interest Groups, Rent-Seeking \& Lobbying & Local government finance; centralisation vs local autonomy; accountability \\
9 & Behavioural Economics & Lobbying (IFA); professional associations; planning; regulatory capture \\
10 & Institutions, Development \& Political Regimes & Nudges; voter heuristics; misinformation \\
11--12 & Project 1 Presentations + Discussion & - \\
\bottomrule
\end{tabular}
\end{frame}

\begin{frame}{Weekly Readings (1/3)}
\small
\begin{tabular}{p{1.5cm}p{8.5cm}}
\toprule
\textbf{Week} & \textbf{Reading} \\
\midrule
Week 1 & \textbf{Recommended:} William F.\ Shughart II (2018) -- Public Choice Theory \\
& \textit{Optional:} Buchanan (2003) -- Politics Without Romance \\
\addlinespace
Week 2 & \textbf{Recommended:} Feddersen (2004) -- Paradox of Not Voting \\
& \textit{Optional:} Blais (2006) -- What Affects Voter Turnout? \\
\addlinespace
Week 3 & \textbf{Recommended:} Holcombe (1989) -- The Median Voter Model in Public Choice Theory \\
& \textit{Optional:} Lee et al.\ (2004) -- Do Voters Affect or Elect Policies? \\
\bottomrule
\end{tabular}
\end{frame}


\begin{frame}{Weekly Readings (2/3)}
\small
\begin{tabular}{p{1.5cm}p{8.5cm}}
\toprule
\textbf{Week} & \textbf{Reading} \\
\midrule
Week 4 & \textbf{Recommended:} Shi \& Svensson (2006) -- Political Budget Cycles \\
& \textbf{Recommended:} Alesina \& Summers (1993) -- Central Bank Independence \\
& \textit{Optional:} Schuknecht (2000) -- Fiscal Policy Cycles \\
\addlinespace
Week 5 & \textbf{Recommended:} Wyckoff (1990) -- The Simple Analytics of Slack-Maximising Bureaucracy \\
& \textit{Optional:} Dunleavy (1985) -- Bureaucrats and State Growth \\
\addlinespace
Week 6 & \textbf{Recommended:} Oates (1999) -- An Essay on Fiscal Federalism \\
& \textit{Optional:} Tiebout (1956) -- A Pure Theory of Local Expenditures \\
\bottomrule
\end{tabular}
\end{frame}


\begin{frame}{Weekly Readings (3/3)}
\small
\begin{tabular}{p{1.5cm}p{8.5cm}}
\toprule
\textbf{Week} & \textbf{Reading} \\
\midrule
\textbf{Week 7} & \textbf{READING WEEK - No assigned readings} \\
\addlinespace
Week 8 & \textbf{Recommended:} Olson (1965) -- \textit{The Logic of Collective Action}, Chapter 1 \\
& \textit{Optional:} Tullock (1967) -- The Welfare Costs of Tariffs \\
& \textit{Optional:} Stigler (1971) -- The Theory of Economic Regulation \\
\addlinespace
Week 9 & \textbf{Recommended:} Thaler \& Sunstein (2003) -- Libertarian Paternalism \\
& \textit{Optional:} DellaVigna (2009) -- Psychology and Economics \\
\addlinespace
Week 10 & \textbf{Recommended:} Acemoglu et al.\ (2005) -- Institutions and Growth \\
& \textbf{Recommended:} Rodrik et al.\ (2004) -- Institutions Rule \\
& \textit{Optional (self-study):} Mueller Chapter 18 -- Dictatorship \& Revolution \\
& \textit{Optional (self-study):} Gandhi \& Przeworski (2006) -- History of Dictators \\
\bottomrule
\end{tabular}
\end{frame}

\begin{frame}{Phase 1: The Demand Side (Voters)}
\textbf{Week 2: The Paradox of Voting}
\begin{itemize}
\item Why vote if your vote doesn't matter? ($P \cdot B - C$)
\item Rational, expressive, and ethical voter models
\item Evidence from turnout studies
\end{itemize}


\textbf{Week 3: The Median Voter Theorem}
\begin{itemize}
\item Single-peaked preferences and policy convergence
\item Why parties look similar
\item When does the theorem break down?
\end{itemize}


\textit{Key question:} If democracy responds to voter preferences, why do voters participate and what do they want?
\end{frame}

\begin{frame}{Phase 2: The Supply Side (Politicians \& Bureaucrats)}
\textbf{Week 4: Political Competition and Macroeconomic Performance}
\begin{itemize}
\item Do politicians manipulate the economy for votes?
\item Political business cycles, opportunistic vs.\ partisan politics
\item Time inconsistency and inflation bias
\item Central bank independence as institutional solution
\item Evidence from elections worldwide
\end{itemize}

\vspace{0.5cm}

\textbf{Week 5: Bureaucracy}
\begin{itemize}
\item Budget-maximising bureaucrats (Niskanen)
\item Information asymmetries and agency problems
\item Public vs.\ private provision of services
\end{itemize}

\vspace{0.5cm}

\textit{Key question:} How do politicians and bureaucrats respond to voter demands?
\end{frame}


\begin{frame}{Phase 3: Interest Groups \& Structures}
\textbf{Week 6: Federalism \& Decentralization}
\begin{itemize}
\item Vertical structure of government
\item Tiebout competition and "voting with your feet"
\item Intergovernmental grants and the flypaper effect
\end{itemize}


\textbf{Week 8: Interest Groups, Rent-Seeking \& Lobbying}
\begin{itemize}
\item Olson: why small groups organise and large groups free-ride
\item Tullock: welfare losses from competing for transfers
\item Stigler: regulation as capture (who benefits from regulation?)
\item Concentrated benefits vs.\ dispersed costs
\end{itemize}


\textit{Key question:} Who shapes policy when political influence is unevenly organised?
\end{frame}

\begin{frame}{Phase 4: Critique \& Synthesis}
\textbf{Week 9: Behavioural Economics}
\begin{itemize}
\item Biases, heuristics, bounded rationality
\item Nudges and choice architecture
\item When do behavioural frictions change political outcomes?
\end{itemize}


\textbf{Week 10: Institutions, Development \& Political Regimes}
\begin{itemize}
\item What makes some countries rich and others poor?
\item Rule of law, property rights, political stability
\item The long-run effects of institutions and path dependence
\item How do institutions differ in democracies vs.\ autocracies?
\item \textit{Self-study materials on dictatorship \& revolution available on Loop}
\end{itemize}


\textit{Key question:} Which institutions sustain growth and accountability over time, and across different political regimes?
\end{frame}

\section{Assessment}

\begin{frame}{Grading Breakdown}
\begin{tabular}{lc}
\toprule
\textbf{Assessment} & \textbf{Weight} \\
\midrule
Project 1: Paper Critique \& Presentation & 45\% \\
Project 2: Grading LLM & 30\% \\
Irish Context Presentation (groups) & 25\% \\
\bottomrule
\end{tabular}

\vspace{0.5cm}
\textbf{Philosophy:} Three substantial assessments that reward genuine learning and application.


\small{No traditional exam - forget memorisation, this is about understanding and applying theory.}
\end{frame}

\begin{frame}{Key Deadlines: Mark Your Calendars}
\begin{tabular}{p{4.5cm}p{5.5cm}}
\toprule
\textbf{Assignment} & \textbf{Deadline} \\
\midrule
Group Formation (Irish Context) & End of Week 2 (Friday, 23 Jan 2026) \\
\addlinespace
Project 1: Paper Declaration & End of Week 4 (Friday, 6 Feb 2026) \\
\addlinespace
Project 1: Written Submission & End of Week 6 (Friday, 20 Feb 2026) \\
\addlinespace
\textbf{Week 7: READING WEEK} & \textbf{23--27 Feb 2026 - No class} \\
\addlinespace
Project 2: Topics Assigned & Start of Week 8 (Monday, 2 Mar 2026) \\
\addlinespace
Project 2: Submission & End of Week 12 (Friday, 3 Apr 2026) \\
\addlinespace
Irish Context Presentations & 48 hours before your assigned week \\
\bottomrule
\end{tabular}

\small{\textbf{Note:} Use reading week to prepare your Project 1 presentation.}
\end{frame}

\begin{frame}{Project 1: Paper Critique (45\%)}
\textbf{The Task:}
\begin{itemize}
\item Find an interesting Public Choice paper published in 2024/2025 or "forthcoming"
\item Write a critical review (50\% of Project 1 = 22.5\% of final grade)
\item Present to class in Weeks 11--12 (50\% of Project 1 = 22.5\% of final grade)
\end{itemize}

\textbf{The Irish Connection:}
\begin{itemize}
\item Your review must include: "How would this apply (or not apply) to Irish politics?" Tip: use your politics/law expertise
\end{itemize}

\textbf{Key Dates:}
\begin{itemize}
\item Declare paper choice: \textbf{End of Week 4 (Friday, 6 Feb 2026)} (email me)
\item Submit review: \textbf{End of Week 6 (Friday, 20 Feb 2026)} (via Loop)
\item Present: Weeks 11--12 (7 minutes + questions)
\end{itemize}
\end{frame}

\begin{frame}{Project 1: What I'm Looking For}
\textbf{Your critique should address:}
\begin{enumerate}
\item What is the research question?
\item What is the main argument/finding?
\item What is the methodology? (Theoretical model? Empirical analysis?)
\item \textbf{Critical evaluation:} What are the strengths and weaknesses?
\begin{itemize}
\item Is the theory convincing?
\item Is the evidence robust?
\item What are the limitations?
\end{itemize}
\item \textbf{The "So what?":} Why does this matter? What should we study next?
\item \textbf{Irish application:} Would this work in Ireland? Why or why not?
\end{enumerate}


\small{Length: 2000--2500 words (strict limit). Quality over quantity.}
\end{frame}


\begin{frame}{Project 1: Presentation Rubric}
\textbf{Your presentation is worth 50\% of Project 1 (= 22.5\% of final grade)}

\vspace{0.2cm}
\textbf{Grading criteria:}
\begin{itemize}
\item \textbf{Content clarity (20\%):} Can you explain the paper's argument clearly?
\item \textbf{Critical analysis (25\%):} Do you demonstrate genuine understanding of strengths/weaknesses?
\item \textbf{Irish application (20\%):} Is your Irish connection thoughtful and specific?
\item \textbf{Delivery (15\%):} Clear speaking, good pacing, effective use of slides?
\item \textbf{Q\&A responses (20\%):} Can you answer questions about the paper and your analysis?
\end{itemize}

\vspace{0.2cm}
\textbf{Remember:} 7 minutes + questions. Rehearse. Both understanding and communication count.
\end{frame}

\begin{frame}{Project 1: Finding Papers}
\textbf{Where to look:}
\begin{itemize}
\item \textit{Public Choice} (the journal)
\item \textit{Journal of Public Economics}
\item \textit{American Political Science Review}
\item \textit{American Economic Review}
\item \textit{Quarterly Journal of Economics}
\item Working papers: NBER, CEPR, IZA
\end{itemize}


\textbf{Topics that count as "Public Choice":}
\begin{itemize}
\item Voting behaviour, electoral systems, political competition
\item Bureaucracy, government spending, taxation
\item Interest groups, lobbying, corruption
\item Institutions and development
\item Anything applying economic methods to political questions
\end{itemize}


\textbf{First come, first served on paper selection.}
\end{frame}

\begin{frame}{Project 2: Grading LLM (30\%)}
\textbf{Part 1: Craft Your Question (40\% of Project 2)}
\begin{itemize}
\item I'll randomly assign you a topic from the module \textbf{(email Week 8: Monday, 2 March 2026)}
\item Write a question for an LLM (e.g., ChatGPT, Claude, Gemini)
\item Write a 500-word motivation explaining:
\begin{itemize}
\item Why this question is important
\item How it relates to class material (be specific)
\item What gap it fills in our knowledge
\item Why the answer would be useful for policy or economics
\end{itemize}
\end{itemize}

\textbf{Part 2: Grade the LLM's Output (60\% of Project 2)}
\begin{itemize}
\item Generate the essay using your question
\item Assess: accuracy, relevance, referencing, depth of argument
\item Identify Irish context errors
\item Provide detailed written feedback explaining your grade
\end{itemize}

\textbf{Submission:} \textbf{End of Week 12 (Friday, 3 April 2026)} via Loop
\end{frame}

\begin{frame}{Project 2: Why Are We Doing This?}
\textbf{The Reality:}
\begin{itemize}
\item LLMs such as GPT, Claude, Gemini, LLaMA, Mistral, Qwen, and DeepSeek are somewhat good at writing essays
\item You'll encounter AI in your careers
\item Learning to \textit{evaluate} AI output is a "crucial" (haha) skill
\end{itemize}


\textbf{What you'll (hopefully) learn:}
\begin{itemize}
\item How to ask good questions (harder than it sounds)
\item What makes an argument strong vs.\ weak
\item How to identify errors, omissions, and shallow reasoning
\item What standards academic work should meet
\item Critical thinking about sources and evidence
\end{itemize}

\textbf{The Irish twist:}
\begin{itemize}
\item You must identify where the LLM gets Irish context wrong
\item Example: Does it understand PR-STV? Coalition formation? Irish institutions?
\end{itemize}
\end{frame}

\begin{frame}{Project 2: What Makes a Good Question?}
\textbf{Your question should:}
\begin{itemize}
\item Build on class material: "In class we discussed X and Y, but not how Z affects..."
\item Fill a knowledge gap: "The theory predicts A, but what about B in the Irish context?"
\item Have practical relevance: "Understanding this would help policymakers/economists because..."
\item Be specific rather than generic
\end{itemize}
\end{frame}

\begin{frame}{Project 2: What Makes a Good Question?}

\textbf{Bad examples:}
\begin{itemize}
\item "Explain the median voter theorem" (too basic, just asks for definition)
\item "What is rent-seeking?" (no connection to class discussion)
\end{itemize}


\textbf{Good examples:}
\begin{itemize}
\item "How does Ireland's PR-STV system affect the median voter theorem's predictions about policy convergence?"
\item "What role does rent-seeking by professional associations play in Irish healthcare waiting lists?"
\end{itemize}
\end{frame}

\begin{frame}{Irish Context Presentations (25\%)}
\textbf{The Format:}
\begin{itemize}
\item Work in \textbf{pairs or trios}
\item Present for 20 minutes at start of class (Weeks 3--6, 8--10)
\item Apply \textit{last week's} theory to Irish politics
\end{itemize}

\textbf{Your presentation should:}
\begin{enumerate}
\item Briefly recap the theory (5 min)
\item Present specific Irish examples/evidence (10 min)
\item Pose 2--3 discussion questions for the class (5 min)
\end{enumerate}

\textbf{Submission requirements:}
\begin{itemize}
\item Slides + 1-page summary due 48 hours before your presentation
\item Submit via Loop assignment dropbox
\item All group members must be present and contribute
\end{itemize}
\end{frame}

\begin{frame}{Irish Context Presentations: Why Groups?}
\textbf{You teach me, I teach you:}
\begin{itemize}
\item You know Irish politics/law better than I do
\item I know the economic theory
\item Together, we build understanding
\end{itemize}

\textbf{Working in groups:}
\begin{itemize}
\item Makes it less scary
\item Allows division of labour
\item Builds accountability
\item Reflects real-world collaboration
\item Trios can cover more ground with deeper analysis
\end{itemize}

\textbf{What I'm looking for:}
\begin{itemize}
\item Accurate understanding of the theory
\item Specific, concrete Irish examples (not vague generalities)
\item Critical thinking: Does the theory fit? Where does it break down?
\item Thoughtful discussion questions
\end{itemize}
\end{frame}


\begin{frame}{Irish Context Presentations: Grading Rubric}
\textbf{Irish Context presentation = 25\% of final grade}

\small
\begin{tabular}{p{3.5cm}p{1.2cm}p{5.5cm}}
\toprule
\textbf{Criterion} & \textbf{Weight} & \textbf{What I'm Looking For} \\
\midrule
Understanding of Theory & 20\% & Accurate recap of last week's theory; clear explanation \\
\addlinespace
Irish Examples & 30\% & Specific examples with evidence; not vague generalities \\
\addlinespace
Critical Analysis & 25\% & Does theory fit Irish reality? Where does it work/break? \\
\addlinespace
Discussion Questions & 15\% & Thoughtful questions that engage the class \\
\addlinespace
Presentation Quality & 10\% & Clear delivery, good slides, time management, equal contribution \\
\bottomrule
\end{tabular}
\end{frame}

\begin{frame}{Academic Integrity: Let's Be Clear}
\textbf{On AI and LLMs:}
\begin{itemize}
\item You may use LLMs to \textit{understand} concepts and brainstorm
\item You may NOT submit AI-generated text as your own work
\item All written work must be in your own words
\item Project 2 is literally about evaluating AI - I'll know if you used it improperly
\end{itemize}


\textbf{On collaboration:}
\begin{itemize}
\item Irish Context Presentations: Collaboration is required (you're in groups)
\item Projects 1 \& 2: Individual work only
\item Discussing ideas is fine; copying text is not
\end{itemize}


\end{frame}

\section{Expectations \& Advice}

\begin{frame}{What I Expect From You}
\begin{itemize}
\item \textbf{Attend class:} We only meet 2 hours/week. Miss at your peril.
\item \textbf{Do the readings:} They'll help you understand the theory and prepare for assessments.
\item \textbf{Engage actively:} Ask questions, challenge ideas, bring examples.
\item \textbf{Meet deadlines:} Late submissions have penalties.
\item \textbf{Be professional:} Presentations are important. Rehearse. Show up prepared.
\item \textbf{Teach me about Ireland:} I'm relying on your expertise.
\end{itemize}


\textbf{Office hours:} Please use them, especially before your presentations.
\end{frame}

\begin{frame}{What You Can Expect From Me}
\begin{itemize}
\item \textbf{Clear explanations:} I'll break down complex theory step-by-step.
\item \textbf{Real-world applications:} Theory without examples is useless.
\item \textbf{Timely feedback:} I'll grade fairly and return work promptly.
\item \textbf{Accessibility:} Email me, attend office hours, I'm here to help.
\item \textbf{Respect for your knowledge:} You bring politics/law expertise I don't have.
\item \textbf{High standards:} I'll push you because this material is important.
\end{itemize}


\textbf{My goal:} You leave this course able to analyse political events like an economist.
\end{frame}

\begin{frame}{How to Succeed in This Course}
\begin{enumerate}
\item \textbf{Don't fall behind:} Each week builds on the previous. Catch up immediately if you miss class.
\item \textbf{Think like an economist:} Ask "What are the incentives?" for every political phenomenon.
\item \textbf{Connect theory to reality:} Read the news with Public Choice glasses on.
\item \textbf{Start projects early:} Don't wait until the deadline.
\item \textbf{Use your interdisciplinary background:} Economics + politics + law = powerful combination.
\item \textbf{Embrace debate:} Public Choice is controversial. Challenge ideas (respectfully).
\end{enumerate}


\small{This is not a memorisation course. Understanding $>$ recitation.}
\end{frame}

\begin{frame}{Recommended Resources}
\textbf{Textbook (optional but helpful):}
\begin{itemize}
\item Mueller, D.\ (2003). \textit{Public Choice III}. Cambridge University Press.
\end{itemize}


\textbf{Accessible introductions:}
\begin{itemize}
\item Buchanan, J.\ \& Tullock, G.\ (1962). \textit{The Calculus of Consent}.
\item Tullock, G.\ (1967). "The Welfare Costs of Tariffs, Monopolies, and Theft."
\item Olson, M.\ (1965). \textit{The Logic of Collective Action}.
\end{itemize}


\textbf{For current applications:}
\begin{itemize}
\item \textit{Public Choice} journal
\item VoxEU / VoxDev (short research summaries)
\item NBER working papers
\end{itemize}

\end{frame}

\section{Why Study Public Choice?}

\begin{frame}{Why Does This Matter?}
\textbf{A concrete promise:} By the end of this course, you will have an economic theory for why housing policy is difficult to change, hospital waiting lists persist, and some taxes are politically untouchable.

\textbf{For your careers:}
\begin{itemize}
\item \textbf{Policy analysts:} Understand why policies fail (not just what policies should be)
\item \textbf{Political advisors:} Predict how institutions shape behaviour
\item \textbf{Lawyers:} Appreciate the economic logic behind legal rules
\item \textbf{Business:} Navigate regulation and government interaction
\item \textbf{Advocacy/NGOs:} Design effective campaigns and coalitions
\end{itemize}


\textbf{For citizenship:}
\begin{itemize}
\item Be a more informed voter
\item Understand why government is hard (not just blame politicians)
\item Recognise tradeoffs and unintended consequences
\end{itemize}
\end{frame}

\begin{frame}{The Big Questions}
By the end of this course, you should be able to answer:

\begin{enumerate}
\item Why do people vote (or not vote)?
\item Why do democracies sometimes produce inefficient outcomes?
\item Why do bureaucracies grow larger than optimal?
\item Why do small interest groups win over the majority?
\item Why do some countries become rich while others stay poor?
\item When does competition improve government, and when does it fail?
\item How should we design political institutions?
\end{enumerate}


\textbf{Bonus:} You'll understand Irish politics much better.
\end{frame}

\begin{frame}{A Warning \& An Invitation}
\textbf{Warning:}
\begin{itemize}
\item Public Choice can be cynical. "Politics without romance."
\item You might become disillusioned with democracy.
\item But understanding problems is the first step to solving them.
\end{itemize}

\pause
\vspace{0.5cm}
\textbf{Invitation:}
\begin{itemize}
\item This material is \textit{challenging} but rewarding.
\item You'll see the world differently after this course.
\item Engage, debate, push back - that's how we learn.
\end{itemize}

\pause
\vspace{0.5cm}
\begin{center}
\textbf{Let's make this a great semester!}
\end{center}
\end{frame}

\section{Next Steps}

\begin{frame}{Before Next Week}
\textbf{Action items:}
\begin{enumerate}
\item Check Loop + Github - all materials will be posted there and here https://github.com/BeaGietner/PublicChoice
\item Read the syllabus carefully (especially assessment details)
\item \textbf{Form groups (pairs or trios) for Irish Context Presentations}
\begin{itemize}
\item Email me your group names by \textbf{end of Week 2 (Friday, 23 January 2026)}
\item I'll assign you a week
\end{itemize}
\item Begin thinking about Project 1 paper selection
\item Read introductory material on Public Choice (posted on Loop)
\end{enumerate}

\vspace{0.5cm}
\textbf{Next week:} The Paradox of Voting - Why do rational people vote?
\end{frame}

\begin{frame}{Questions?}
\begin{center}
\Large Open discussion time!

\vspace{1cm}
Ask me anything about:
\begin{itemize}
\item The course structure
\item Assessments
\item Public Choice in general
\item My background / research
\item Anything else
\end{itemize}
\end{center}
\end{frame}

\begin{frame}{Final Thought}
\begin{center}
\large "The man of system... is apt to be very wise in his own conceit... he seems to imagine that he can arrange the different members of a great society with as much ease as the hand arranges the different pieces upon a chess-board." \\
\vspace{0.5cm}
- Adam Smith, \textit{The Theory of Moral Sentiments}
\end{center}

\vspace{0.5cm}
\pause
\begin{center}
Public Choice reminds us: \\
\textbf{The pieces on the chessboard have their own agendas.}
\end{center}
\end{frame}

\section{Irish Context Presentations: Guidance}

\begin{frame}{Week 3 Presentation}
\textbf{Applies Week 2: Paradox of Voting}

\textbf{Theme:} Turnout, participation, abstention

\textbf{Example prompts:}
\begin{itemize}
\item Why is voter turnout in Irish general elections relatively high compared to many OECD countries? Does rational choice explain this, or do we need expressive/ethical voting models?
\item Compare turnout in Irish referendums vs.\ general elections. What does Public Choice predict, and does the data fit?
\item Does compulsory voting (hypothetical) make sense for Ireland given the paradox of voting?
\end{itemize}

\textbf{What you're testing:} Rational ignorance, expressive voting, civic duty, institutional context.
\end{frame}

\begin{frame}{Week 4 Presentation}
\textbf{Applies Week 3: Median Voter Theorem}

\textbf{Theme:} Party positioning, PR-STV, convergence

\textbf{Example prompts:}
\begin{itemize}
\item Do Fianna Fáil, Fine Gael, and Labour converge toward the median Irish voter? On which policy dimensions do they diverge, and why?
\item How does PR-STV change the predictions of the median voter theorem compared to two-party systems?
\item Is Sinn Féin's rise consistent with median voter logic or a challenge to it?
\end{itemize}

\textbf{What you're testing:} Single-peaked preferences, multidimensionality, institutional constraints.
\end{frame}

\begin{frame}{Week 5 Presentation}
\textbf{Applies Week 4: Political Competition and Macroeconomic Performance}

\textbf{Theme:} Fiscal cycles, budgets, opportunism, central bank independence

\textbf{Example prompts:}
\begin{itemize}
\item Is there evidence of opportunistic fiscal policy in Irish budgets before elections?
\item Do Irish governments manipulate taxation, spending, or public sector pay in election years?
\item How does ECB membership constrain Irish fiscal policy? Is this constraint beneficial from a Public Choice perspective?
\item Does Ireland's lack of monetary policy independence prevent political business cycles, or do they manifest in other ways?
\end{itemize}

\textbf{What you're testing:} Opportunistic vs.\ partisan models, timing of policy, voter myopia, institutional constraints on policy manipulation.
\end{frame}

\begin{frame}{Week 6 Presentation}
\textbf{Applies Week 5: Bureaucracy}

\textbf{Theme:} HSE, civil service, agency problems

\textbf{Example prompts:}
\begin{itemize}
\item Does the HSE behave like a budget-maximising bureaucrat? Use budget data and service outcomes to assess.
\item Are Irish public sector agencies characterised by slack, overproduction, or inefficiency?
\item How do information asymmetries between ministers and civil servants shape policy outcomes?
\end{itemize}

\textbf{What you're testing:} Niskanen model, principal--agent problems, information asymmetry.
\end{frame}

\begin{frame}{Week 7: READING WEEK}
\begin{center}
\Large \textbf{No Class - No Irish Context Presentation}

\vspace{1cm}

\textbf{Use this week to:}
\begin{itemize}
\item Prepare your Project 1 presentation slides
\item Practice your 7-minute presentation
\item Review your paper one more time
\item Get feedback from peers or come to office hours
\end{itemize}
\end{center}
\end{frame}

\begin{frame}{Week 8 Presentation}
\textbf{Applies Week 6: Federalism \& Decentralisation}

\textbf{Theme:} Centralisation, local government, intergovernmental finance

\textbf{Example prompts:}
\begin{itemize}
\item Is Ireland "too centralised"? What would Public Choice predict about efficiency and accountability under greater decentralisation?
\item Do Irish local authorities have meaningful fiscal autonomy? What are the incentive effects?
\item Does the Tiebout model make any sense in the Irish context? Why or why not?
\item How do intergovernmental grants affect local government behaviour in Ireland?
\end{itemize}

\textbf{What you're testing:} Tiebout sorting, fiscal federalism, accountability, mobility, centralisation vs.\ local autonomy.
\end{frame}

\begin{frame}{Week 9 Presentation}
\textbf{Applies Week 8: Interest Groups, Rent-Seeking \& Lobbying}

\textbf{Theme:} Lobbying, organised interests, concentrated benefits

\textbf{Example prompts:}
\begin{itemize}
\item Analyse the IFA (Irish Farmers' Association) through Olson's Logic of Collective Action. Why are they effective?
\item Do professional associations (Law Society, medical bodies) engage in rent-seeking in Ireland? Where do you see it?
\item How do planning regulations create opportunities for rent-seeking or regulatory capture?
\item Why do consumers remain rationally ignorant about policies that cost them money?
\end{itemize}

\textbf{What you're testing:} Free-rider problem, concentrated vs.\ dispersed costs, capture, rent-seeking, collective action.
\end{frame}

\begin{frame}{Week 10 Presentation}
\textbf{Applies Week 9: Behavioural Economics}

\textbf{Theme:} Nudges, heuristics, misinformation

\textbf{Example prompts:}
\begin{itemize}
\item Evaluate an Irish behavioural intervention (e.g.\ tax compliance, health, energy). Does it work? Why?
\item Do Irish voters exhibit framing effects or status quo bias in referendums? (Consider recent referendum campaigns.)
\item How does misinformation affect political outcomes in Ireland? Is this predictable under behavioural models?
\item What role do heuristics (mental shortcuts) play in Irish voting behaviour?
\end{itemize}

\textbf{What you're testing:} Bounded rationality, heuristics, behavioural biases in politics, nudges.
\end{frame}

\begin{frame}{Weeks 11--12: Project 1 Presentations}
\textbf{These weeks are reserved for Project 1 presentations.}

\textbf{Format:}
\begin{itemize}
\item 7 minutes presentation + questions
\item Submit slides in advance (48 hours before your assigned week)
\item Focus: research question, method, main finding, and your critique
\end{itemize}

\textbf{Reminder:} Your critique must include an Irish connection (does the argument travel to Ireland? why or why not?).
\end{frame}

\end{document}